\chapter{Tesztelés és minőségbiztosítás}

A WaccApp frontendjének tesztelése a fejlesztési folyamat fontos visszacsatolási szakasza volt. A rendszer esetében nem volt elegendő azt ellenőrizni, hogy az egyes oldalak megnyílnak-e, vagy hogy a fő gombok technikailag elindítanak-e egy műveletet. Egy WhatsApp Business API-ra épülő kommunikációs felületnél a használhatóság ugyanúgy része a minőségnek, mint a funkcionális működés. A felhasználónak értenie kell, milyen adatot kell megadnia, mikor indítható el egy küldés, mit jelent egy hibaüzenet, és hol tudja később visszakeresni az elküldött vagy beérkezett tartalmakat.

A tesztelés ezért elsősorban a WaccApp tényleges felhasználói folyamataira épült. Az ellenőrzések középpontjában az \path{index.html} küldőfelülete, a \path{conv.html} beszélgetésnézete, a \path{messages.html} és a \path{sent-messages.html} listaoldalai, a \path{questionnaire.html} és a \path{template.html} kérdőíves felületei, valamint a \path{filter.html} szűrési nézete állt. Ezek az oldalak a korábbi fejezetekben külön modulokként jelentek meg, a tesztelés során azonban az volt a lényeges, hogy együtt is működőképes kommunikációs folyamatot alkossanak.

A minőségbiztosítás során nemcsak a sikeres eseteket kellett kipróbálni, hanem azokat a helyzeteket is, amelyek a valós használatban bizonytalanságot okozhatnak. Ilyen lehet a hibás telefonszám, az üres üzenetmező, a hiányzó sablonparaméter, a nem megfelelő fájl, a hibás kérdőívkapcsolat, a rosszul megadott dátumintervallum vagy egy olyan szűrés, amely nem ad találatot. Ezek a hibák nem mindig látványos rendszerleállásként jelennek meg. Sokszor éppen az a probléma, hogy a felület működik, de a felhasználó nem kap elég egyértelmű választ arra, hogy mi történt. A tesztelés egyik célja ezért az volt, hogy ezek a bizonytalan helyzetek felismerhetők és javíthatók legyenek.

A tesztelés közben többször előfordult, hogy egy funkció technikailag működött, használat közben mégis bizonytalannak érződött. Ilyen helyzet volt például, amikor egy művelet elindult, de a felület nem adott elég gyors visszajelzést, vagy amikor egy üres találati lista mellett nem volt egyértelmű, hogy valóban nincs adat, vagy a lekérdezés nem futott le. Ezek a tapasztalatok megerősítették, hogy a frontend tesztelésénél nemcsak a sikeres API-választ, hanem a felhasználó által érzékelt folyamatot is vizsgálni kell.

\section{A tesztelés célja és szempontrendszere}

A tesztelés elsődleges célja annak igazolása volt, hogy a WaccApp fő frontendfolyamatai a gyakorlatban is végigvihetők. Ez nemcsak azt jelentette, hogy az egyes funkciók külön-külön elérhetők, hanem azt is, hogy a felhasználó egy teljes kommunikációs munkafolyamatot el tud végezni. Egy üzenet vagy kérdőív előkészítése, elküldése, állapotának értelmezése, majd későbbi visszakeresése csak együtt ad teljes képet a rendszer működéséről.

A tesztelési szempontrendszer első része a funkcionális működés ellenőrzése volt. Az \path{index.html} oldalon azt kellett kipróbálni, hogy elindíthatók-e az egyedi szöveges üzenetek, a sablonos üzenetek, a kérdőívek, a médiafájlok és az időzített műveletek. A \path{messages.html}, a \path{sent-messages.html} és a \path{conv.html} esetében az volt a kérdés, hogy a kommunikációs adatok áttekinthetők-e, elkülönülnek-e a bejövő és kimenő tartalmak, valamint egy adott kontakt beszélgetése értelmezhető formában jelenik-e meg. A \path{questionnaire.html} és a \path{template.html} tesztelésekor a kérdések, válaszok, sablonok és következő lépések kezelése került előtérbe. A \path{filter.html} ellenőrzése pedig arra irányult, hogy a megadott feltételek alapján valóban értelmezhető eredmények jelenjenek meg.

A második szempont a hibás vagy hiányos adatbevitel kezelése volt. A WaccApp nem tekinthető megbízható frontendnek, ha csak ideális kitöltés mellett működik. Ezért külön figyelmet kaptak azok az esetek, amikor a felhasználó hibás telefonszámot ad meg, üresen hagyja az üzenetmezőt, nem tölti ki a sablonparamétereket, fájl nélkül próbál médiát küldeni, vagy időzített műveletnél hiányzó, illetve múltbeli időpontot választ. Ezeknél a teszteknél nem az volt a cél, hogy a frontend minden hibát automatikusan kijavítson, hanem az, hogy a hiba oka érthetően jelenjen meg, és a felhasználó tudja, mit kell módosítania.

A harmadik szempont a használhatóság és a reszponzív működés ellenőrzése volt. A WaccApp több különálló HTML-nézetből épül fel, ezért fontos volt megvizsgálni, hogy a felület kisebb kijelzőn is használható marad-e. Az \path{index.html} esetében a küldési mezők és gombok kezelhetősége, a \path{conv.html} esetében a beszélgetés görgethetősége és olvashatósága, a \path{filter.html} esetében pedig a több szűrési mező áttekinthetősége volt meghatározó. A kérdőívszerkesztő oldalaknál külön figyelmet igényelt, hogy a kérdések, válaszok és kapcsolódások kisebb képernyőn is követhető sorrendben jelenjenek meg.

A negyedik szempont a regressziós ellenőrzés volt. A WaccApp fejlesztése során több funkció egymásra épült, ezért egy módosítás könnyen hatással lehetett más nézetekre is. Egy kérdőíves logikai javítás például érinthette az \path{index.html} kérdőívindítását, a \path{conv.html} beszélgetésmegjelenítését vagy a \path{filter.html} kérdőíves szűrését. Emiatt egy-egy javítás után nem volt elegendő csak az adott rész gyors kipróbálása, hanem a kapcsolódó folyamatokat is újra ellenőrizni kellett.

A manuális tesztelés során a cél az volt, hogy a WaccApp legfontosabb frontendfolyamatai ellenőrizhető módon kipróbálhatók legyenek. Emiatt a tesztelés nem kizárólag egyes gombok vagy mezők működésére irányult, hanem teljesebb felhasználói folyamatokra: üzenetküldésre, sablonhasználatra, kérdőívkezelésre, szűrésre, médiamegjelenítésre és időzített műveletekre. A részletes manuális tesztforgatókönyvek a 15.3. mellékletben találhatók. A manuális tesztelés fő ellenőrzési területeit pedig a \ref{tab:manualis-teszteles}. táblázat foglalja össze.

\begin{table}[H]
\centering
\caption{A manuális tesztelés fő ellenőrzési területei.}
\label{tab:manualis-teszteles}
\vspace{0.6em}
\footnotesize
\setlength{\tabcolsep}{3pt}
\renewcommand{\arraystretch}{1.25}

\begin{tabularx}{\textwidth}{|>{\raggedright\arraybackslash}p{0.18\textwidth}|>{\raggedright\arraybackslash}p{0.23\textwidth}|Y|>{\raggedright\arraybackslash}p{0.14\textwidth}|}
\hline
\textbf{Tesztterület} &
\textbf{Érintett frontendnézetek} &
\textbf{Ellenőrzés célja} &
\textbf{Eredmény} \\
\hline

Üzenetküldési folyamatok &
\path{index.html}, \path{sent-messages.html}, \path{conv.html} &
Annak ellenőrzése, hogy az egyedi és sablonos üzenetküldés elindítható-e, a visszajelzés megjelenik-e, és az elküldött tartalom később visszakereshető-e. &
Megfelelt \\
\hline

Hibás vagy hiányos adatbevitel &
\path{index.html}, \path{questionnaire.html}, \path{template.html}, \path{filter.html} &
Annak vizsgálata, hogy a frontend felismeri-e a hiányzó címzettet, üres üzenetet, hiányzó sablonparamétert, hibás időpontot vagy félrevezető szűrési feltételt. &
Megfelelt \\
\hline

Kérdőíves működés &
\path{questionnaire.html}, \path{template.html}, \path{index.html}, \path{filter.html} &
Annak ellenőrzése, hogy a kérdőívek létrehozhatók, indíthatók, válaszfüggő logikával követhetők, majd később szűrhetők és egyszerűen összesíthetők-e. &
Megfelelt \\
\hline

Kommunikáció visszakövetése &
\path{messages.html}, \path{sent-messages.html}, \path{conv.html} &
Annak vizsgálata, hogy a beérkező és elküldött üzenetek listában, illetve beszélgetésnézetben is értelmezhetően jelennek-e meg. &
Megfelelt \\
\hline

Média- és fájlkezelés &
\path{index.html}, \path{conv.html}, public/uploads, public/sent\_media &
Annak ellenőrzése, hogy a médiafájl küldése elindítható-e, és a korábban elküldött kép vagy dokumentum a beszélgetési kontextusban értelmezhetően jelenik-e meg. &
Megfelelt \\
\hline

Szűrés és válaszösszesítés &
\path{filter.html} &
Annak vizsgálata, hogy a név, telefonszám, dátum, üzenettípus, kérdőív és válasz alapján történő keresés, valamint az egyszerű darabszám szerinti összesítés működik-e. &
Megfelelt \\
\hline

Reszponzív és használhatósági ellenőrzés &
Fő frontendnézetek &
Annak ellenőrzése, hogy a felületek kisebb kijelzőn is kezelhetők, olvashatók és követhetők maradnak-e. &
Megfelelt; kisebb fejlesztési lehetőségekkel \\
\hline

\end{tabularx}
\end{table}

A táblázat alapján látható, hogy a tesztelés a fő frontendfolyamatok gyakorlati használhatóságát vizsgálta. A sikeres esetek mellett hibás vagy hiányos adatbeviteli helyzetek is szerepeltek, mert ezek mutatják meg, hogy a frontend validációs és visszajelzési logikája valóban támogatja-e a felhasználót. A tesztek eredménye alapján a WaccApp fő funkciói manuális ellenőrzéssel végigvihetők voltak, ugyanakkor az is láthatóvá vált, hogy a későbbi fejlesztések során az automatizált tesztelés és a részletesebb naplózás tovább erősíthetné a rendszer megbízhatóságát.

\section{Üzenetküldési és sablonküldési tesztek}

Az üzenetküldési tesztek az \path{index.html} oldalon kezdődtek, mert ez a WaccApp központi küldőfelülete. Az egyszerű szöveges üzenet tesztelésekor a felhasználói folyamat a címzett megadásából, az üzenetszöveg kitöltéséből és a küldés indításából állt. A teszt akkor volt sikeres, ha a rendszer nemcsak elküldte a kérést a háttérrendszer felé, hanem közben egyértelműen jelezte a folyamat állapotát, majd a művelet eredményét is megjelenítette. Ez azért volt fontos, mert visszajelzés nélkül a felhasználó könnyen újra megnyomhatja a küldés gombot, mivel nem lehet biztos abban, hogy a kérés valóban elindult.

Az üzenetküldés ellenőrzése nem zárult le a küldés gomb megnyomásával. A sikeres küldés után azt is meg kellett vizsgálni, hogy az elküldött tartalom megjelenik-e a \path{sent-messages.html} oldalon, illetve visszakövethető-e a \path{conv.html} beszélgetésnézetben. Ez a lépés azért volt lényeges, mert a WaccApp nem egyszerű küldőfelületként működik. A rendszer célja az is, hogy a kommunikáció később értelmezhető maradjon, ezért az elküldött üzeneteknek a kommunikációs előzmények részeként is meg kellett jelenniük.

A hibás telefonszám tesztelésekor azt kellett ellenőrizni, hogy a frontend felismeri-e a nyilvánvaló formai hibákat. Ha a telefonszám hiányos, rossz szerkezetű vagy nem illeszkedik a várt formátumhoz, akkor a felületnek már a küldés előtt jeleznie kell a problémát. A frontend ebben az esetben nem végleges hitelesítést végez, mert azt a backend és a külső platform válaszai határozzák meg. A kliensoldali ellenőrzés szerepe inkább az, hogy a legegyszerűbb hibák ne jussanak el feleslegesen a szerverig, és a felhasználó azonnal javítani tudjon.

Az üres üzenet küldésének tesztje szintén az alapműködéshez tartozott. Ilyenkor a felhasználó megadja a címzettet, de nem ír üzenetszöveget. A megfelelő működés az, ha a rendszer nem indítja el a küldést, hanem jelzi, hogy az üzenettartalom megadása szükséges. Ez egyszerű teszthelyzetnek tűnik, mégis fontos, mert egy üres üzenet elküldése használati szempontból értelmetlen művelet lenne, és rontaná a rendszer megbízhatóságát.

A sablonküldés tesztelésekor a sablon kiválasztása, a címzett megadása és a szükséges paraméterek kitöltése került előtérbe. A sikeres teszt feltétele az volt, hogy a frontend a kiválasztott sablonhoz illeszkedő adatokat továbbítsa, és a küldés után egyértelmű visszajelzést adjon. A hiányzó sablonparaméter külön hibatesztként jelent meg. Ilyenkor a küldés nem indulhat el úgy, mintha minden adat rendelkezésre állna, mert egy hiányos sablon később backendoldali vagy platformszintű hibát okozhatna. A felületnek ezért már az előkészítés során segítenie kell abban, hogy a felhasználó észrevegye a hiányzó adatot.

Az üzenetküldési és sablonküldési tesztek egyik tanulsága az volt, hogy a sikeres API-hívás önmagában nem elég a jó felhasználói élményhez. A felhasználó számára az is fontos, hogy a küldés előtti ellenőrzés, a folyamat közbeni állapotjelzés és a későbbi visszakereshetőség együtt működjön. Ha ezek közül bármelyik hiányzik, a rendszer technikailag akár működhet is, de a használat mégis bizonytalanná válhat.

\section{Kérdőív- és sablonfolyamatok tesztelése}

A kérdőíves működés tesztelése a WaccApp egyik legkényesebb minőségbiztosítási területe volt, mert itt nem egyszerű adatküldésről, hanem több lépésből álló kommunikációs logikáról van szó. A \path{questionnaire.html} oldalon először azt kellett ellenőrizni, hogy a felhasználó létre tud-e hozni egy szöveges vagy számalapú válaszokra épülő kérdőívet. A teszt során kérdést kellett rögzíteni, válaszlehetőségeket kellett megadni, majd be kellett állítani, hogy egy adott válasz után melyik következő kérdés jelenjen meg. A sikeres működés feltétele az volt, hogy a kérdőív szerkezete menthető és később újra felhasználható legyen.

A kérdőíves tesztelésnél nem volt elegendő egyetlen kérdés létrehozását ellenőrizni. A valódi próba az volt, hogy a kérdésekből összeállított lánc végigjárható marad-e. Ha egy válasz rossz következő elemre mutat, vagy hiányzik a folytatás, akkor a kérdőív nem látványos hibával áll meg, hanem egyszerűen nem folytatódik. Ez frontend szempontból azért problémás, mert a felhasználó ilyenkor nem tudja eldönteni, hogy ő hagyott ki valamit, vagy a rendszer nem működik megfelelően. Emiatt a kérdőívláncok ellenőrzése a tesztelés egyik központi része lett.

A kérdőívindítás tesztje az \path{index.html} és a kérdőíves adatstruktúra kapcsolatát vizsgálta. A felhasználónak ki kellett választania egy korábban létrehozott kérdőívet, majd el kellett indítania azt egy megadott címzettre. A teszt akkor tekinthető megfelelőnek, ha a rendszer valóban a kiválasztott kérdőív kezdő kérdését indította el. Ez a pont azért lényeges, mert ha a frontend rossz kérdőívazonosítót, rossz témát vagy hiányos adatot küld tovább, akkor a folyamat már az első lépésnél félrecsúszik.

A válasz utáni következő kérdés ellenőrzése külön tesztforgatókönyvet igényelt. A kérdőív akkor működik valóban, ha a beérkező válasz nemcsak rögzül, hanem a folyamat ennek alapján tovább is halad. A teszt során azt kellett vizsgálni, hogy egy válasz után a megfelelő következő kérdés indul-e el, illetve más válasz esetén másik útvonalra lép-e a rendszer. Ez a működés különbözteti meg a WaccApp kérdőívkezelését egy egyszerű statikus kérdéslistától.

A hibás kérdőívlogika tesztelésekor olyan helyzeteket kellett kipróbálni, amikor egy válaszhoz nem tartozik következő kérdés, vagy a kérdőívben hiányzik egy hivatkozott elem. Ilyenkor az volt a cél, hogy a rendszer ne viselkedjen megtévesztően. Ideális esetben a frontend már szerkesztés közben jelzi, hogy a kérdőív szerkezete hiányos. Ha a hiba csak futás közben derül ki, akkor legalább annak kell világosnak lennie, hogy a folyamat nem folytatható megfelelően. Ez a teszt azért volt fontos, mert a kérdőívlogika megbízhatósága nem látványkérdés, hanem a teljes kommunikációs folyamat feltétele.

A sablonos kérdőívek tesztelése a \path{template.html}, a \path{template-menu.html} és a \path{template-single.html} oldalakhoz kapcsolódott. Ezeknél azt kellett ellenőrizni, hogy a gombos sablonfolyamatok létrehozhatók-e, a sablonnevek és válaszlehetőségek következetesen kezelhetők-e, valamint a gombos válasz után a megfelelő következő sablon következik-e. A sablonos kérdőív tesztelése azért különbözik a szöveges kérdőívtől, mert itt a folyamat erősebben kötődik a WhatsApp Business Platform által elvárt sablonlogikához. Emiatt a sablonparaméterek, gombszövegek és sablonláncok ellenőrzése külön figyelmet igényelt. [R1], [R2], [R3]

A kérdőíves tesztek tapasztalata alapján ez a modul csak akkor tekinthető stabilnak, ha a szerkesztés, az indítás, a válaszfeldolgozás és a visszakeresés együtt működik. Nem elég, ha a kérdőív létrehozható, de nem indítható el. Az sem elegendő, ha elindítható, de a válasz után nem a megfelelő kérdés következik. A WaccApp kérdőíves funkciójának minőségét ezért a teljes kérdéslánc kipróbálása alapján lehetett megítélni.

\section{Szűrési és időzített műveletek ellenőrzése}

A \path{filter.html} tesztelésekor a cél az volt, hogy a rendszer ne csak adatokat listázzon, hanem a felhasználó által megadott feltételek alapján értelmezhető eredményt adjon. A név és telefonszám szerinti szűrésnél azt kellett ellenőrizni, hogy a rendszer megfelelően kapcsolja-e össze a kontaktadatokat, és nem enged-e félrevezető párosítást. Ez azért volt lényeges, mert egy rosszul összerendelt név és telefonszám hibás következtetésekhez vezethetne a kommunikáció visszakeresésekor.

A dátumintervallumos szűrés tesztelésekor a felhasználó kezdő és záró dátumot adott meg, majd ellenőrizni kellett, hogy az eredmények valóban az adott időszakhoz tartoznak-e. Hibás dátumlogika esetén, például ha a kezdődátum későbbi, mint a záródátum, a frontendnek ezt nem szabad csendben elfogadnia. Egy ilyen lekérdezés technikailag lefuthatna, de a felhasználó számára értelmetlen vagy félrevezető eredményt adna. A tesztelés ezért nemcsak a lekérdezés lefutását, hanem az eredmény értelmezhetőségét is vizsgálta.

A kérdőívhez kapcsolódó szűrésnél azt kellett kipróbálni, hogy a felhasználó képes-e egy adott kérdőívhez, kérdéshez vagy válaszhoz tartozó adatokat visszakeresni. Ez a funkció azért fontos, mert a WaccApp-ban a kérdőív nem önmagában áll, hanem később értelmezhető adatforrást is jelent. Ha a válaszok nem lennének szűrhetők, akkor a kérdőívkezelés gyakorlati értéke jelentősen csökkenne. A \path{filter.html} ezért a tesztelés során nem egyszerű keresőoldalként, hanem a kommunikációs és kérdőíves adatok értelmező felületeként jelent meg.

A beszélgetésnézet tesztelése a \path{conv.html} oldalon történt. A teszt célja az volt, hogy egy kiválasztott kontakt kommunikációs előzményei ne szétszórt rekordokként, hanem időrendi beszélgetésként jelenjenek meg. Ellenőrizni kellett, hogy a bejövő és kimenő üzenetek elkülönülnek-e, a sablonos üzenetek értelmezhető tartalomként jelennek-e meg, és a beszélgetés olvasása közben a felhasználó követni tudja-e a kommunikáció menetét. Ez a nézet különösen fontos, mert a WaccApp egyik leginkább felhasználóközeli része.

A média- és fájlkezelés tesztelése az \path{index.html}, a \path{conv.html}, valamint a public/uploads és public/sent\_media mappákhoz kapcsolódó működés alapján történt. A sikeres teszt során a felhasználó kiválasztott egy támogatott képet vagy dokumentumot, a rendszer megjelenítette a kiválasztott fájlhoz kapcsolódó információt, majd a küldés után a tartalom később is visszakövethető maradt. A hibás fájlkiválasztás tesztelésekor nem támogatott fájltípust vagy hiányzó fájlt kellett kezelni. Ilyenkor a megfelelő működés az volt, ha a rendszer nem indította el félrevezető módon a küldést, hanem jelezte, hogy a fájl nem alkalmas a műveletre.

A beszélgetésnézet tesztelésekor különbséget jelentett, hogy a médiatartalom egyszerű hivatkozásként vagy ténylegesen értelmezhető vizuális elemként jelent meg. Egy fájlelérési útvonal technikailag információt ad, de a felhasználó számára kevésbé beszédes. A \path{conv.html} esetében ezért fontos tesztszempont volt, hogy az elküldött képek a beszélgetés részeként jelenjenek meg, ne csak olyan adatként, amelyet külön meg kell nyitni. Ez a megoldás közelebb viszi a felületet ahhoz az üzenetküldési élményhez, amelyet a felhasználók a mindennapi kommunikációs alkalmazásokból már ismernek. [R13], [R14]

Az időzített küldések tesztelése azért kapott külön figyelmet, mert ezeknél a művelet rögzítése és végrehajtása időben elválik egymástól. Az \path{index.html} felületén ellenőrizni kellett, hogy a felhasználó be tud-e állítani jövőbeli időpontot szöveges üzenethez, sablonhoz, kérdőívhez vagy médiatartalomhoz. A sikeres teszt nemcsak azt jelentette, hogy az időpont mentésre került, hanem azt is, hogy a rendszer a művelethez tartozó kötelező adatokat is ellenőrizte. Ilyen adat lehet a címzett, az üzenettartalom, a sablonparaméter, a kérdőív kiválasztása vagy a csatolt fájl.

A hibás időzítési helyzetek tesztelésekor több esetet kellett kipróbálni. Ilyen volt a hiányzó időpont, a múltbeli időpont, a hiányos üzenettartalommal beállított időzítés, valamint a médiafájl nélküli médiaütemezés. Ezekben az esetekben a frontendnek már a rögzítés előtt jeleznie kellett, hogy a művelet nem teljes. Ez különösen fontos, mert időzített küldésnél a hiba később derülne ki, amikor a felhasználó már nem feltétlenül figyeli aktívan a rendszert. A tesztelés ezért azt is vizsgálta, hogy a várakozó, sikeres és sikertelen állapotok később értelmezhetően jelennek-e meg.

\section{Reszponzív ellenőrzések, javítások és minőségbiztosítási tanulságok}

A reszponzív tesztelés során a WaccApp fő nézeteit eltérő kijelzőméreteken kellett ellenőrizni. Az \path{index.html} esetében azt kellett vizsgálni, hogy a küldési módok, beviteli mezők és gombok mobilnézetben is használhatók maradnak-e. A \path{conv.html} oldalon külön figyelmet kapott a beszélgetési terület görgetése, a fejléc és az üzenetküldő sáv helyfoglalása, valamint az üzenetek olvashatósága. A \path{filter.html} esetében a több mezőből álló szűrőfelület kezelhetőségét kellett ellenőrizni kisebb kijelzőn. A kérdőívszerkesztő oldalaknál pedig az volt a fő kérdés, hogy a kérdésblokkok, válaszlehetőségek és kapcsolódások mobilon is követhető sorrendben jelennek-e meg. [R6], [R7], [R9]

A tesztelés során több olyan helyzet is láthatóvá vált, amikor a felület technikailag működött, de a felhasználó számára nem adott elég egyértelmű visszajelzést. Ilyen lehetett, amikor egy művelet után nem derült ki azonnal, hogy a küldés folyamatban van-e, vagy amikor egy üres találati lista nem jelezte világosan, hogy valóban lefutott a szűrés. Ezek nem klasszikus programhibák, mégis rontják a rendszer használhatóságát, mert bizonytalanságot keltenek. A frontend minősége ilyen esetekben azon múlik, hogy a felület képes-e érthetően közölni, mi történt.

A hibák több csoportba sorolhatók. Az egyik csoportot a beviteli hibák alkották, például a hibás telefonszám, az üres üzenetmező, a hiányzó sablonparaméter vagy a hiányos kérdőívkapcsolat. A másik csoportba a megjelenítési hibák tartoztak, például amikor egy beszélgetésben a médiatartalom nem volt elég jól értelmezhető, vagy a szűrési eredmények nem adtak elég világos visszajelzést. A harmadik csoportot az állapotkezelési problémák jelentették, amikor a felület nem kommunikálta elég egyértelműen, hogy egy művelet folyamatban van, sikerült vagy hibára futott.

A javítások során több esetben nem új funkciót kellett hozzáadni, hanem a meglévő működést kellett érthetőbbé tenni. Ilyen volt a hiányzó mezők pontosabb jelzése, a médiatartalmak beszélgetésben történő megjelenítésének javítása, a szűrési eredmények egyértelműbb kezelése, valamint az időzített műveletek státuszainak követhetőbb megjelenítése. A kérdőíves folyamatoknál a logikai kapcsolatok ellenőrzése vált fontosabbá, mert egy hibás következő lépés a teljes kérdésláncot megakaszthatta.

A tesztelés egyik legfontosabb tanulsága az volt, hogy a WaccApp frontendje nem kezelhető elszigetelt oldalak gyűjteményeként. Egy funkció értéke sokszor csak több nézet együttműködésén keresztül látható. Egy üzenet elküldése például az \path{index.html} oldalon indul, de a \path{sent-messages.html} és a \path{conv.html} nézetekben válik visszakövethetővé. Egy kérdőív a \path{questionnaire.html} vagy a \path{template.html} oldalon készül elő, az \path{index.html} oldalon indítható, a válaszok pedig később a beszélgetésnézetben és a \path{filter.html} szűrési felületén értelmezhetők.

A fejezet alapján a WaccApp tesztelése nem pusztán lezáró ellenőrzés volt, hanem a frontendfejlesztési döntések gyakorlati igazolása is. A tesztelés megmutatta, hogy a rendszer használhatósága nem egyetlen látványos megoldáson múlik, hanem a kisebb döntések következetességén. Ilyen döntés a hiányzó adatok jelzése, a küldési állapot megjelenítése, a beszélgetések időrendi kezelése, a médiaelemek értelmezhető megjelenítése, a kérdőívláncok ellenőrzése és a szűrési eredmények egyértelmű visszajelzése. Ezek együtt biztosítják, hogy a WaccApp frontendje ne csak technikailag működjön, hanem a gyakorlatban is kiszámíthatóan használható legyen. [R6], [R7], [R13], [R14] 
